% !TEX program = tectonic
\documentclass[9pt,a4paper]{extarticle}
\usepackage[top=0.8cm,bottom=0.8cm,left=0.9cm,right=0.9cm]{geometry}
\usepackage[T1]{fontenc}\usepackage[utf8]{inputenc}\usepackage{lmodern}
\usepackage[french]{babel}
\usepackage{amsmath,amssymb}\usepackage{textcomp}\usepackage{xcolor}\usepackage{booktabs}\usepackage{array}
\newcolumntype{R}[1]{>{\raggedright\arraybackslash}p{#1}}
\usepackage{enumitem}\usepackage{mdframed}\usepackage{multicol}\usepackage{graphicx}
\graphicspath{{figs/}}
\definecolor{bleu}{RGB}{20,77,144}\definecolor{vert}{RGB}{20,110,60}
\definecolor{rouge}{RGB}{160,30,30}\definecolor{grisf}{RGB}{110,110,110}
\newcommand{\fig}[2]{\begin{center}\includegraphics[width=\linewidth]{#1}\end{center}\vspace{-2.5mm}{\scriptsize\color{grisf} #2}\par\vspace{0.3mm}}
\pagestyle{empty}\setlength{\parindent}{0pt}\setlength{\parskip}{1.1pt}
\setlength{\abovedisplayskip}{2.5pt}\setlength{\belowdisplayskip}{2.5pt}
\setlength{\abovedisplayshortskip}{1pt}\setlength{\belowdisplayshortskip}{1pt}
\renewcommand{\arraystretch}{0.95}
\setlength{\columnsep}{6mm}\setlength{\columnseprule}{0.3pt}
\renewcommand{\columnseprulecolor}{\color{grisf!40}}
\newcommand{\et}[2]{\vspace{1.3mm}{\color{bleu}\bfseries #1 $\cdot$ #2}\par\vspace{-0.2mm}}
\newcommand{\res}[1]{\vspace{0.3mm}{\small\color{bleu}$\Rightarrow$\ #1}\par\vspace{0.2mm}}
\newcommand{\obs}[1]{\vspace{0.2mm}\hangindent=3.4mm\hangafter=0{\color{bleu}\footnotesize\textbullet}~{\small #1}\par\vspace{0.2mm}}
\newcommand{\lire}[1]{{\scriptsize\color{vert}$\blacktriangleright$~\textbf{lire :} #1}\par}
\newenvironment{pts}{\begin{itemize}[leftmargin=3.4mm,itemsep=0.2mm,topsep=0.4mm,parsep=0pt,label={\color{bleu}\footnotesize\textbullet}]}{\end{itemize}}
\newmdenv[linecolor=bleu!55,backgroundcolor=bleu!5,linewidth=0.5pt,innertopmargin=2.5pt,
 innerbottommargin=2.5pt,innerleftmargin=3pt,innerrightmargin=3pt,skipabove=1.5pt,skipbelow=1.5pt]{boite}
\begin{document}

\begin{center}
{\large\bfseries Stratégie par ticker --- suite de tests}
\end{center}
\vspace{-1.5mm}

\begin{multicols}{2}

% ═══════════ É1 ═══════════
\et{É1}{La donnée --- le périmètre}
\begin{pts}
\item On part de la table déjà propre (pipeline S1/S2) --- aucun nettoyage refait.
\item On ne garde qu'un sous-ensemble : actions cotées, 2014--2026, série de prix exploitable.
\end{pts}
\begin{center}\scriptsize
\begin{tabular}{@{}l r l@{}}
\toprule
table clean & 134\,464 & \\
$\to$ actions seules & 128\,974 & {\color{grisf}$-$5\,490 : 2\,190 ETF, 3\,300 inconnu} \\
$\to$ années 2014--2026 & 128\,269 & {\color{grisf}$-$705 : donnée trop mince} \\
$\to$ prix exploitable & \textbf{113\,369} & {\color{grisf}$-$14\,900} \\
\bottomrule
\end{tabular}
\end{center}
\lire{chaque ligne retire un morceau ; la colonne = nombre de transactions qui restent.}
{\scriptsize\color{grisf} 2\,755 titres $\cdot$ 350 membres $\cdot$ 57\,445 achats, 55\,924 ventes.\par}

% ═══════════ É2 ═══════════
\et{É2}{Dater les lots (FIFO) et purger le dormant}
\begin{pts}
\item Chaque achat empile un lot : (date, nombre de parts).
\item Une vente consomme les lots \textbf{les plus anciens} d'abord.
\item L'âge ne choisit pas le lot vendu, il l'\textbf{étiquette} : \og récent\fg{} si détenu $\le 2$ ans, \og vieux\fg{} sinon.
\item Le signal ne garde que la part \textbf{récente} d'une vente (le dormant ne signale pas un délit d'initié).
\end{pts}
\[ \gamma_k = \frac{\sum_{\ell \subset k} \mathbf{1}\{h_\ell \le \bar a\}\, q_\ell}{\sum_{\ell \subset k} q_\ell} \in [0,1] \]
\begin{pts}
\item $s_k$ : jour de bourse où le trade $k$ s'exécute.
\item $q_k = A_k / P(s_k)$ : nombre de parts ($A_k$ = montant, $P$ = prix).
\item $h_\ell$ : durée de détention du lot $\ell$ vendu, en jours de bourse.
\item $\bar a = 504$ j : seuil \og vieux\fg{} = 2 ans de bourse.
\item $\gamma_k$ : part de la vente $k$ portant sur du récent ($0$ = tout vieux, $1$ = tout récent).
\item $\mathbf{1}\{\cdot\}$ = 1 si la condition est vraie, 0 sinon $\cdot$ $\ell \subset k$ : lot $\ell$ (de $q_\ell$ parts) consommé par la vente $k$.
\end{pts}
\fig{durees.png}{}
\lire{$x$ = durée de détention des ventes (jours) ; barre rouge = $\bar a$ (2 ans) ; à droite = lots dormants, retirés du signal.}
\begin{center}\scriptsize
\begin{tabular}{@{}l r@{}}
\toprule
$\gamma{=}1$ vente entièrement récente (gardée) & 79,9 \% \\
$\gamma{=}0$ entièrement vieille (retirée) & \textbf{17,1 \%} \\
$0{<}\gamma{<}1$ partielle & 3,1 \% \\
\midrule
masse vendue gardée dans le signal & \textbf{89,5 \%} des 2,11 Md\$ \\
\bottomrule
\end{tabular}
\end{center}
\lire{les 3 premières lignes = quelle fraction des ventes tombe dans ce cas ; dernière = part du \$ vendu qui alimente le signal.}

% ═══════════ É3 ═══════════
\et{É3}{Le signal ticker et la pondération}
\begin{pts}
\item Sur une fenêtre de $W$ jours, on somme les achats et la part récente des ventes de chaque titre.
\item On retient un titre s'il est net-acheté par \textbf{au moins deux membres}.
\item Chaque membre pèse pareil ($1/K_t$), qu'il ait mis 11 M\$ ou 15 k\$, et répartit son poids entre ses titres.
\end{pts}
\[ NetBuy_{i,t} = Buy_{i,t} - Sell_{i,t}\cdot\gamma \]
\[ S_t = \{\, i : NetBuy_{i,t} > 0 \;\wedge\; nBuyers_{i,t} \ge 2 \,\} \]
\[ w_{i,t} = \frac{1}{K_t} \sum_{m \in \mathcal{M}_t} \frac{\mathbf{1}\{ i \in B_{m,t} \}}{|B_{m,t}|} \]
\begin{pts}
\item $Buy,\ Sell$ : montants achetés / vendus du titre $i$ dans la fenêtre.
\item $nBuyers$ : nombre de membres distincts ayant acheté $i$.
\item $S_t$ : les titres retenus le mois $t$ ; $B_{m,t}$ : ceux de $S_t$ achetés par le membre $m$.
\item $\mathcal{M}_t$ : les membres ayant acheté $\ge1$ titre de $S_t$ (membres actifs) ; $K_t = |\mathcal{M}_t|$ ; $w_{i,t}$ : poids du titre $i$ ($\sum_i w_i = 1$).
\end{pts}

% ═══════════ É4 ═══════════
\et{É4}{Choisir la fenêtre $W$ (test de concentration)}
\begin{pts}
\item On teste 4 fenêtres et on tranche sur la \textbf{concentration}, jamais sur le rendement.
\end{pts}
\[ W^\star = \min\bigl\{\, W : \mathbb{P}[\max_i w_{i,t} > 10\,\%] \le 20\,\% \,\bigr\} \]
\begin{pts}
\item $\max_i w_{i,t}$ : poids de la plus grosse ligne, un mois donné (règle : $>10\%$ dans $\le20\%$ des mois).
\end{pts}
\begin{center}\scriptsize
\begin{tabular}{@{}l r r r r@{}}
\toprule
$W$ & 21 j & \textbf{42 j} & 63 j & 126 j \\
\midrule
$|S_t|$ médian & 28 & \textbf{64} & 91 & 156 \\
$K_t$ médian & 19 & \textbf{32} & 39 & 53 \\
$\max_i w_i$ médian & 11,0 \% & \textbf{7,2 \%} & 6,3 \% & 5,4 \% \\
mois $> 10\,\%$ & 58 \% & \textbf{17 \%} & 9 \% & 5 \% \\
\bottomrule
\end{tabular}
\end{center}
\lire{colonnes = les 4 fenêtres testées ; $|S_t|$ = titres retenus/mois ; $K_t$ = membres actifs/mois ; $\max w$ = plus grosse ligne ; dernière ligne = fréquence de sur-concentration (le critère).}
\res{$W = 42$ j retenu.}
\obs{à $W{=}21$ une ligne montait à 50 \% et $K_t$ tombait à 2 ; à $W{=}42$, $|S_t|$ médian 64, $K_t$ 32, jamais vide.}

% ═══════════ encadré lecture ═══════════
\begin{boite}
{\footnotesize\bfseries\color{bleu}Comment lire les résultats} {\scriptsize\color{grisf}--- notations.}
\begin{pts}
\item \textbf{excès/an} : $\displaystyle x_Y = \prod_{t\in Y}(1{+}r_t) - \prod_{t\in Y}(1{+}r_{\mathrm{SPY},t})$ pour l'année $Y$, puis moyenne sur les années.
\item \textbf{$t$} : $t = \bar x \,/\, (\sigma_x/\sqrt{n})$, $n$ = nombre d'années ($|t|\ge2$ : significatif à 5 \%).
\item \textbf{$\alpha,\ \beta$} : régression $r_t - r_{f,t} = \alpha + \beta\,(r_{\mathrm{SPY},t} - r_{f,t}) + \varepsilon_t$, annualisé $\alpha_{\mathrm{an}} = 252\,\alpha$ ($r_f$ = taux sans risque, $\varepsilon$ = résidu).
\item \textbf{NAV} : $\mathrm{NAV}_T = 100\prod_{t\le T}(1{+}r_t)$, base 100 en 2014 (SPY : 517).
\end{pts}
\end{boite}

% ═══════════ É5 ═══════════
\et{É5}{Test 1 --- la base}
\begin{pts}
\item Long seul (on achète, jamais de vente à découvert).
\item Fenêtre du signal sur la \textbf{transaction} $\tau$ (date où le membre a agi).
\item Aucune contrainte de rotation. Coupe en fin de mois, exécution au \textbf{close J+1} (cours de clôture du lendemain).
\end{pts}
\res{excès $\mathbf{+1{,}71\,\%}$/an $\cdot$ $t\,1{,}03$ $\cdot$ $\alpha\,+1{,}05\,\%$ $\cdot$ $\beta\,1{,}02$ $\cdot$ NAV \textbf{622}}
\obs{7/13 années positives ; $\beta \approx 1$ ; NAV 622 $>$ SPY 517.}
\fig{beta_glissant.png}{}
\lire{courbe bleue = $\beta$ de la base sur 1 an glissant ; il reste entre \textbf{0,93 et 1,11} (médiane 1,02) ; pointillés = $\beta{=}1$ (marché) et $\beta$ plein-échantillon 1,02.}

% ═══════════ É6 ═══════════
\et{É6}{Test 2 --- + cap de rotation 25 \%/mois}
\begin{pts}
\item On ne réalloue pas d'un coup : on avance \emph{partiellement} vers la cible, d'au plus 25 \%/mois.
\end{pts}
\[ \mathrm{TO} = \tfrac12 \textstyle\sum_i |w^{\text{cible}}_i - w^{\text{avant}}_i| \]
\[ \lambda = \min\!\bigl(1,\ \tau_{\max}/\mathrm{TO}\bigr), \qquad \tau_{\max}{=}25\,\% \]
\begin{pts}
\item $\mathrm{TO}$ : rotation voulue ce mois ($0$ = rien ne bouge, $1$ = tout change).
\item $\lambda$ : fraction du chemin vers la cible qu'on parcourt réellement.
\end{pts}
\res{excès $\mathbf{+0{,}27\,\%}$/an $\cdot$ $t\,0{,}23$ $\cdot$ $\alpha\,-0{,}16\,\%$ $\cdot$ $\beta\,1{,}01$ $\cdot$ NAV \textbf{531}}
\obs{le cap ramène la rotation de 59 \% à 25 \%/mois ; NAV 531 (base 622) ; excès $+0{,}27\,\%$/an ($t\,0{,}23$) contre $+1{,}71\,\%$ ($t\,1{,}03$).}

% ═══════════ É7 ═══════════
\et{É7}{Test 3 --- + jambe short (long-short)}
\begin{pts}
\item Au lieu d'ignorer les titres net-\textbf{vendus}, on les shorte (on parie sur leur baisse).
\item On ne short que ceux corroborés par $\ge2$ vendeurs de stock récent.
\end{pts}
\[ r = r^{+} - \theta\, r^{-}, \qquad \theta = 1 \]
\begin{pts}
\item $S^{-}$ : titres à flux net négatif, $\ge2$ vendeurs récents ; $r^{+},r^{-}$ : rendement du livre long / short.
\item $\theta$ : intensité du short ($\theta{=}0$ redonne le long seul).
\end{pts}
\res{excès $\mathbf{-12{,}8\,\%}$/an $\cdot$ $t\,-3{,}19$ $\cdot$ $\beta\,-0{,}02$ $\cdot$ NAV \textbf{122}}
\obs{$|S^{-}|$ médian 45 ; les titres net-vendus \emph{montent} après la vente ; $\beta$ passe à $\sim$0.}
\fig{extensions.png}{}
\lire{$y$ = NAV base 100 (échelle log) ; A = base (T1), B = $+$cap 25 \% (colle à A), C = long-short (s'effondre), gris = SPY.}

% ═══════════ É8 short-only ═══════════
\et{É8}{Test 4 --- short-only (la jambe short seule)}
\begin{pts}
\item On enlève toute position longue : on \textbf{shorte} seulement les titres net-vendus ($\ge2$ vendeurs récents).
\end{pts}
\res{excès $\mathbf{-26{,}6\,\%}$/an $\cdot$ $t\,-3{,}40$ $\cdot$ $\alpha\,-2{,}36\,\%$ $\cdot$ $\beta\,\mathbf{-1{,}03}$ $\cdot$ NAV \textbf{16}}
\obs{les titres net-vendus montent : shortés seuls, 100~\$ tombent à 16 ; $\beta\,{\approx}\,{-}1$ = à l'inverse du marché.}

% ═══════════ É9 δ ═══════════
\et{É9}{Test 5 --- à la date de divulgation $\delta$}
\begin{pts}
\item Même base (long seul), mais on agit à la \textbf{divulgation} $\delta$ (quand le trade devient public) au lieu de $\tau$.
\item Le FIFO et $\gamma$ restent sur $\tau$ : la durée de détention est une réalité, pas une date de publication.
\end{pts}
\res{excès $\mathbf{+0{,}35\,\%}$/an $\cdot$ $t\,0{,}18$ $\cdot$ $\alpha\,-0{,}03\,\%$ $\cdot$ $\beta\,1{,}01$ $\cdot$ NAV \textbf{530}}
\obs{lag médian de divulgation 27 j ; l'écart cumulé sur SPY passe de $+105$ pts ($\tau$) à $+13$ pts ($\delta$).}
\begin{center}\scriptsize
\begin{tabular}{@{}l r r@{}}
\toprule
excès/an (NAV) & sans cap & cap 25 \% \\
\midrule
$\tau$ (transaction) & $+1{,}71\,\%$ (622) & $+0{,}27\,\%$ (531) \\
$\delta$ (divulgation) & $+0{,}35\,\%$ (530) & $-0{,}49\,\%$ (487) \\
\bottomrule
\end{tabular}
\end{center}
{\color{bleu}\footnotesize\textbullet}~{\small le cap réduit l'excès à $\tau$ comme à $\delta$.\par}
\fig{coins.png}{}
\lire{$y$ = NAV base 100 (log) ; $\tau$ = transaction, $\delta$ = divulgation ; trait plein = sans cap, tireté = cap 25 \% ; pointillé = SPY.}
\begin{center}\scriptsize
\begin{tabular}{@{}l r r@{\quad}l r r@{}}
\toprule
année & exc.\,$\tau$ & exc.\,$\delta$ & année & exc.\,$\tau$ & exc.\,$\delta$ \\
\midrule
2014 & $-1{,}3$ & $+2{,}2$ & 2021 & $+2{,}0$ & $\mathbf{+13{,}2}$ \\
2015 & $-0{,}9$ & $+0{,}1$ & 2022 & $-3{,}1$ & $-1{,}6$ \\
2016 & $+9{,}4$ & $+11{,}1$ & 2023 & $\mathbf{+11{,}5}$ & $-3{,}8$ \\
2017 & $+3{,}0$ & $+3{,}7$ & 2024 & $-8{,}7$ & $-9{,}4$ \\
2018 & $+0{,}9$ & $-2{,}4$ & 2025 & $+10{,}9$ & $+3{,}5$ \\
2019 & $-3{,}4$ & $+3{,}0$ & 2026 & $-1{,}6$ & $-5{,}4$ \\
2020 & $+3{,}5$ & $\mathbf{-9{,}6}$ & \textbf{moy.} & $\mathbf{+1{,}7}$ & $\mathbf{+0{,}4}$ \\
\bottomrule
\end{tabular}
\end{center}
\lire{une ligne = une année ; exc.\,$\tau$ et exc.\,$\delta$ = l'excès sur SPY selon qu'on agit à la transaction ou à la divulgation ; moy. = la colonne excès/an du récap.}

% ═══════════ récap ═══════════
\et{}{Les cinq tests côte à côte}
\begin{center}\scriptsize
\begin{tabular}{@{}l r r r r r@{}}
\toprule
test & excès/an & $t$ & $\alpha$/an & $\beta$ & NAV \\
\midrule
T1 base (long, $\tau$) & $+1{,}71\,\%$ & 1,03 & $+1{,}05\,\%$ & 1,02 & 622 \\
T2 $+$ cap 25 \% & $+0{,}27\,\%$ & 0,23 & $-0{,}16\,\%$ & 1,01 & 531 \\
T3 $+$ long-short ($\theta{=}1$) & $-12{,}8\,\%$ & $-3{,}19$ & $-0{,}25\,\%$ & $-0{,}02$ & 122 \\
T4 short-only ($\tau$) & $-26{,}6\,\%$ & $-3{,}40$ & $-2{,}36\,\%$ & $-1{,}03$ & 16 \\
T5 base, à $\delta$ & $+0{,}35\,\%$ & 0,18 & $-0{,}03\,\%$ & 1,01 & 530 \\
\bottomrule
\end{tabular}
\end{center}
\lire{les 5 colonnes = cf. encadré \og Comment lire les résultats\fg{}.}

% ─── É10 · Test 6 · Comités ───
\et{É10}{Test 6 --- comités pertinents}
\textbf{Ce qu'on fait} (le reste est la base à l'identique, \textbf{en long}) :
\begin{pts}
\item \textbf{Le filtre} : on ne garde que les achats d'un membre \textbf{siégeant sur le comité qui régule le secteur du titre} --- Financial Services$\to$Financials, Armed Services$\to$Industriels, Energy \& Commerce$\to$\{Énergie, Santé\}$\dots$ ($18\,948$ trades, $16{,}7\,\%$).
\item On \textbf{achète} $S^{c}_t$ (aucune vente à découvert), une voix par membre restreinte à ces membres, mensuel, close J+1. \emph{Variante grossière} : comité \og clé\fg{} (drapeau, sans le secteur).
\end{pts}
\[ k \text{ retenu} \iff \text{secteur}(i_k) \in \text{secteurs}\bigl(\text{comités}(m_k)\bigr) \]
\[ w^{c}_{i,t} = \frac{1}{K^{c}_t}\sum_{m\in\mathcal{M}^{c}_t}\frac{\mathbf{1}\{i\in B^{c}_{m,t}\}}{|B^{c}_{m,t}|}, \qquad \textstyle\sum_i w^{c}_i = 1 \]
\begin{center}\scriptsize
\begin{tabular}{@{}l r r r r@{}}
\toprule
& exc.\,$\tau$ & $t_\tau$ & exc.\,$\delta$ & $t_\delta$ \\
\midrule
comité-secteur & $\mathbf{+7{,}67}$ & \textbf{1,60} & $-3{,}65$ & $-0{,}99$ \\
comité-clé (drapeau) & $+1{,}74$ & 0,70 & $-1{,}05$ & $-0{,}56$ \\
\bottomrule
\end{tabular}
\end{center}
\fig{comites.png}{}
\textbf{Ce qu'on observe} :
\begin{pts}
\item À la transaction $\tau$, restreindre aux comités \textbf{double la base} : NAV \textbf{1102} contre 622, $\alpha\,{+}5{,}95\,\%$/an ($t_\alpha$ 1,35).
\item Il \textbf{s'effondre à la divulgation} $\delta$ : NAV 326, $-3{,}65\,\%$/an.
\item \textbf{À nuancer} : seulement $\sim$4 titres/mois --- l'excès annuel part dans tous les sens, et le filtre est très déséquilibré par secteur (ci-dessous) : au seuil usuel, le $+7{,}67$ à $\tau$ ne se distingue pas du hasard.
\end{pts}
\fig{comites_annuel.png}{excès annuel à $\tau$ (bleu) et $\delta$ (vert) : de $-30$ à $+45\,\%$ selon l'année.}
\fig{comites_secteurs.png}{\% des trades de chaque secteur retenus par le filtre : 63 \% Industrie, 0 \% Cons.\,Disc./Immobilier.}

% ═══════════ Produits réels : NANC / GOP / Quiver ═══════════
\et{}{Les produits réels --- ce qu'ils annoncent vs ce qu'on obtient}
Trois produits commercialisent la \og{}copie du Congrès\fg{}. Pour chacun : la règle qu'ils \textbf{disent} suivre, le chiffre qu'ils \textbf{annoncent}, puis \textbf{notre} réplique fidèle de leur règle.

\textbf{1. Quiver Quantitative} (indices \og{}Congress Trading\fg{}). \textbf{Ce qu'ils disent faire} :
\begin{pts}
\item \textbf{Congress Buys} : long les titres \textbf{achetés} par le Congrès, \textbf{pondérés à la taille du trade} (\$), rotation \textbf{hebdo}. \textbf{Long-Short 130/30} : $+$130 \% des achats, $-$30 \% des ventes.
\item On réplique à la transaction $\tau$ (borne haute) \emph{et} à la divulgation $\delta$ (réplicable) ; fenêtre 63 j, rotation mensuelle.
\end{pts}
\[ w^{\text{Buys}}_{i,t} \;=\; \frac{\displaystyle\sum_{k\,:\ achat,\ i_k=i,\ t-W<s_k\le t} A_k}{\displaystyle\sum_{j}\ \sum_{k\,:\ achat,\ i_k=j,\ t-W<s_k\le t} A_k},\qquad \textstyle\sum_i w^{\text{Buys}}_i=1 \]
\[ r^{\text{LS}}_t \;=\; 1{,}3\,r^{\text{long}}_t - 0{,}3\,r^{\text{short}}_t \]
\begin{center}\scriptsize
\begin{tabular}{@{}l r r r@{}}
\toprule
CAGR/an \emph{2020-26} & annoncé & rép.\,$\tau$ & rép.\,$\delta$ \\
\midrule
Congress Buys & 36,2 \% \, $\beta$1,14 & \textbf{22,6 \%} & 21,4 \% \\
Long-Short 130/30 & 34,0 \% \, $\beta$1,15 & \textbf{23,5 \%} & 22,3 \% \\
SPY (même fenêtre) & --- & 21,1 \% & 21,1 \% \\
\bottomrule
\end{tabular}
\end{center}
\begin{center}\scriptsize
\begin{tabular}{@{}l r r r r r@{}}
\toprule
notre réplique (2014-26) & excès/an & $t$ & $\alpha$/an & $\beta$ & NAV \\
\midrule
Congress Buys $\cdot\ \tau$ & $+0{,}84\,\%$ & 0,74 & $-0{,}21\,\%$ & 1,06 & 565 \\
Congress Buys $\cdot\ \delta$ & $-0{,}03\,\%$ & $-0{,}03$ & $-0{,}69\,\%$ & 1,04 & 511 \\
Long-Short $\cdot\ \tau$ & $+1{,}57\,\%$ & 1,13 & $+0{,}32\,\%$ & 1,07 & 606 \\
Long-Short $\cdot\ \delta$ & $+0{,}92\,\%$ & 0,60 & $+0{,}07\,\%$ & 1,04 & 555 \\
\bottomrule
\end{tabular}
\end{center}

\textbf{1bis. Quiver --- version \textit{House uniquement}} (page \og{}U.S. House Long-Short\fg{}, même mécanique, flux restreint à la Chambre). \textbf{Ce qu'ils annoncent} (depuis 2020-04) : \textbf{$+$628,58 \%} \og{}All Time\fg{} $\cdot$ CAGR \textbf{36,94 \%} $\cdot$ $\beta$ 1,09 $\cdot$ Sharpe 1,024 $\cdot$ IR 0,71 $\cdot$ MaxDD $-$24,70 \%. \textbf{Notre réplique fidèle} (House, $\delta$, base 100 au 2020-04) : \textbf{$+$267 \%} $\cdot$ CAGR \textbf{23,2 \%} $\cdot$ $\beta$ 1,04 $\cdot$ Sharpe 0,74 $\cdot$ IR 0,17 $\cdot$ MaxDD $-$31,5 \% \emph{(SPY même départ : $+$230 \%)}. Un contrôle d'identité vérifie qu'en retirant le filtre chambre on retombe exactement sur la réplique Congress ci-dessus.
\begin{center}\scriptsize
\begin{tabular}{@{}l r r r r r r@{}}
\toprule
réplique House (2014-26) & excès/an & $t$ & $\alpha$/an & $\beta$ & Sharpe & NAV \\
\midrule
Congress Buys $\cdot\ \delta$ & $+0{,}61\,\%$ & 0,48 & $-0{,}18\,\%$ & 1,04 & 0,72 & 546 \\
Long-Short $\cdot\ \tau$ & $+1{,}44\,\%$ & 1,09 & $+0{,}28\,\%$ & 1,07 & 0,74 & 600 \\
Long-Short $\cdot\ \delta$ & $+1{,}48\,\%$ & 0,86 & $+0{,}50\,\%$ & 1,04 & 0,74 & 587 \\
Long-Short $\cdot\ \delta\ $ hebdo & $+1{,}68\,\%$ & 1,04 & $+0{,}64\,\%$ & 1,05 & 0,75 & 609 \\
\bottomrule
\end{tabular}
\end{center}
\textbf{Même verdict que la version Congress} : la Chambre ne change rien --- $\approx$ le marché avec du levier, $\alpha\approx 0$. L'écart au 36,94 \% tient à leur boîte noire non divulguée (fenêtre, timing $\tau/\delta$, panier short) et à un panneau de métriques incohérent ($\beta$ 1,09 impose une vol $\ge$ 19 \%/an, incompatible avec leur vol affichée).

\fig{quiver.png}{}

\textbf{2. NANC / GOP} (ETF Subversive) --- \emph{traité à part}
\begin{pts}
\item Ces deux ETF copient les élus d'un parti (Dém.\ $\to$ NANC, Rép.\ $\to$ GOP) et \textbf{publient leur portefeuille chaque jour} : on peut donc comparer notre réplique à la vérité, ligne par ligne.
\item \textbf{Résultat} : on reproduit \textbf{la moitié} du portefeuille NANC et \textbf{moins d'un tiers} du GOP. On retrouve leurs grands noms, pas leurs tailles.
\item \textbf{Pourquoi} : leurs positions valent jusqu'à \textbf{25 fois} le flux déclaré correspondant --- leur fournisseur de données n'est pas les fourchettes publiques, et le gérant ajoute une part de discrétion qu'il admet lui-même.
\item \textbf{Et là non plus, pas d'avantage} : $\beta\approx1$, aucun $t$ au-dessus de 2. NANC bat le S\&P de 2,7 pts/an, GOP le sous-performe de 6,2.
\end{pts}
\lire{le détail (leurs documents, nos trois lectures de leur règle, la composition ligne à ligne, le diagnostic) est dans la fiche dédiée \textbf{FICHE\_NANC\_GOP}, désormais dans \texttt{\_archive/}.}
\end{multicols}

% ═══════════ ANNEXE : extraits sources ═══════════
\vspace{2mm}\hrule\vspace{1.5mm}
{\color{bleu}\bfseries\large Annexe --- extraits sources}\par\vspace{2mm}

\textbf{1. Quiver --- Congress Buys} \emph{(quiverquant.com)} :
{\footnotesize\itshape ``The Congress Buys Strategy tracks the performance of stocks that have been purchased by members of U.S. Congress (or their family). This strategy is weighted based on the reported size of the purchases, with weekly rebalancing.''\par}
\vspace{1.5mm}

\textbf{1. Quiver --- Congress Long-Short} \emph{(quiverquant.com)} :
{\footnotesize\itshape ``The Congress Long-Short Strategy tracks the performance of stocks that have been purchased or sold by members of U.S. Congress (or their family). It takes a long position in stocks that have been purchased, and a short position in stocks which have been sold. This strategy is weighted based on the reported size of the transactions and employs leverage with 130\% long exposure and 30\% short exposure, with weekly rebalancing.''\par}
\vspace{2mm}

\end{document}
